\chapter{Opérations sur les limites et suites géométriques}
\section*{Exercice 10.1 -- Terme dominant}
[Base] Calculer la limite de $(5n^2+2n)/(n^2-3)$ à l'infini.

\section*{Exercice 10.2 -- Conjugué}
[Méthode] Calculer $\lim(\sqrt{n^2+4n}-n)$.

\section*{Exercice 10.3 -- Géométriques}
[Base] Étudier $(0{,}7)^n$, $(-0{,}7)^n$, $(-2)^n$.

\section*{Exercice 10.4 -- Suite affine}
[Méthode] Résoudre $u_0=1$, $u_{n+1}=u_n/2+3$.

\section*{Exercice 10.5 -- Indétermination}
[Critique] Construire deux quotients de suites tendant toutes vers zéro, l'un tendant vers 2 et l'autre vers $+\infty$.

\section*{Exercice 10.6 -- Somme infinie}
[Approfondissement] Définir puis calculer $\sum_{k=1}^{\infty}3^{-k}$.
